\documentclass[9pt,a4paper]{extarticle}
\usepackage{parskip}

% =====================
% Codificação, idioma e layout
% =====================
\usepackage[T1]{fontenc}
\usepackage[utf8]{inputenc}
\usepackage[brazil]{babel}
\usepackage{lmodern}
\usepackage{geometry}
\geometry{margin=2cm}
\usepackage{microtype}

% =====================
% Matemática e símbolos
% =====================
\usepackage{amsmath, amssymb, amsthm, bm, mathtools}

% =====================
% Links
% =====================
\usepackage{hyperref}
\hypersetup{colorlinks=true,linkcolor=black,citecolor=black,urlcolor=blue}

% =====================
% Início do documento
% =====================
\begin{document}

\section*{Roteiro de Aula -- Métodos de Busca Unidirecional}

\subsection*{1. Algoritmo de Descida (visão geral)}

Um \textbf{algoritmo de descida} é um método iterativo para encontrar um ponto que minimiza uma função \( f(\bm{x}) \).
A cada iteração \( k \), escolhe-se:
\begin{itemize}
  \item uma \textbf{direção de descida} \( \bm{d}_k \), tal que \( \nabla f(\bm{x}_k)^\top \bm{d}_k < 0 \);
  \item um \textbf{comprimento de passo} \( \alpha_k > 0 \).
\end{itemize}

A atualização é dada por:
\[
\bm{x}_{k+1} = \bm{x}_k + \alpha_k \bm{d}_k.
\]

O desafio central é determinar um valor apropriado para \( \alpha_k \), o que leva ao \textbf{problema de busca unidirecional}.

\subsection*{2. Problema de Busca Unidirecional}

Dado \( \bm{x}_k \) e uma direção de descida \( \bm{d}_k \), define-se:
\[
\phi(\alpha) = f(\bm{x}_k + \alpha \bm{d}_k), \quad \alpha \ge 0.
\]

O problema consiste em resolver:
\[
\min_{\alpha > 0} \, \phi(\alpha),
\]
ou seja, determinar o valor de \( \alpha \) que reduz mais eficientemente o valor de \( f \) ao longo da direção escolhida.

\subsection*{3. Método de Bracketing}

O \textbf{método de bracketing} busca determinar um intervalo \( [a,b] \) que contenha um mínimo local de \( \phi(\alpha) \), satisfazendo:
\[
\phi(a) > \phi(b) < \phi(c).
\]

\textbf{Etapas:}
\begin{enumerate}
  \item Escolher um ponto inicial \( \alpha_0 \) e um incremento \( \Delta > 0 \);
  \item Avaliar \( \phi(\alpha_0) \) e \( \phi(\alpha_0 + \Delta) \);
  \item Expandir o intervalo enquanto \( \phi \) estiver decrescendo;
  \item Parar quando encontrar um “vale”, isto é, quando a função voltar a crescer.
\end{enumerate}

O bracketing é normalmente usado como etapa inicial antes de métodos mais refinados, como a seção áurea ou interpolação.

\subsection*{4. Método da Seção Áurea}

O \textbf{método da seção áurea} reduz iterativamente o intervalo que contém o mínimo de uma função unimodal, sem necessidade de derivadas.

Define-se a razão áurea:
\[
\tau = \frac{\sqrt{5} - 1}{2} \approx 0{,}618.
\]

\textbf{Etapas:}
\begin{enumerate}
  \item Inicializar um intervalo \( [a,b] \);
  \item Calcular os pontos internos:
  \[
  \alpha_1 = a + (1 - \tau)(b - a), \quad
  \alpha_2 = a + \tau(b - a);
  \]
  \item Avaliar \( \phi(\alpha_1) \) e \( \phi(\alpha_2) \);
  \item Substituir o extremo com valor maior e repetir até que \( |b - a| \) seja pequeno.
\end{enumerate}

Esse método é eficiente, simples e garante convergência quando \( \phi \) é unimodal.

\subsection*{5. Método de Interpolação Quadrática}

O \textbf{método de interpolação quadrática} aproxima \( \phi(\alpha) \) por um polinômio de segunda ordem ajustado a três pontos conhecidos.

\textbf{Etapas:}
\begin{enumerate}
  \item Escolher três pontos distintos \( \alpha_1, \alpha_2, \alpha_3 \);
  \item Ajustar um modelo quadrático:
  \[
  q(\alpha) = a\alpha^2 + b\alpha + c;
  \]
  \item Determinar o ponto que minimiza \( q \):
  \[
  \alpha^* = -\frac{b}{2a};
  \]
  \item Atualizar os pontos, mantendo sempre os que delimitam o mínimo.
\end{enumerate}

Esse método converge rapidamente perto do mínimo, mas pode ser instável se os pontos forem mal espaçados ou o ruído numérico for alto.

\subsection*{6. Busca Inexata -- Regra de Armijo (Backtracking)}

Em muitos algoritmos práticos, não se busca o mínimo exato de \( \phi(\alpha) \), e sim um valor que reduza suficientemente a função.

A \textbf{condição de Armijo} (ou de \emph{decréscimo suficiente}) é:
\[
f(\bm{x}_k + \alpha \bm{d}_k)
\le f(\bm{x}_k) + c_1 \alpha \nabla f(\bm{x}_k)^\top \bm{d}_k,
\]
com \( 0 < c_1 < 1 \) (geralmente \( c_1 = 10^{-4} \)).

\textbf{Algoritmo de Backtracking:}
\begin{enumerate}
  \item Definir \( \alpha \leftarrow \alpha_0 \) (por exemplo, \( \alpha_0 = 1 \));
  \item Enquanto a condição de Armijo não for satisfeita, atualizar:
  \[
  \alpha \leftarrow \rho \alpha, \quad 0 < \rho < 1;
  \]
  \item Retornar o primeiro \( \alpha \) que satisfaz a desigualdade.
\end{enumerate}

Vantagens:
\begin{itemize}
  \item Simples e eficiente;
  \item Evita passos grandes demais;
  \item Muito usado em métodos de Gradiente, Newton e Quasi-Newton.
\end{itemize}

\end{document}
